\bookchapter{Bab 31: Pengenalan Penggambaran Video Pendek dan Penciptaan Kandungan}

\chapterintro{2026, video pendek bukan lagi “mainan orang muda”, ia telah menjadi cara asas semua orang mendapatkan maklumat dan menyatakan diri. Tidak kira anda mahu rekod kehidupan, kongsi pengetahuan, atau buat kerja sampingan, telefon adalah alat penciptaan terbaik anda - tidak perlu beli kamera, tidak perlu belajar Premiere, satu telefon boleh mula. Bab ini bantu anda bermula dari sifar.}

\booksection{Sebelum merakam, fikir dulu: Anda merakam untuk siapa?}

Ramai orang tangkap video kesilapan pertama adalah “angkat telefon terus tangkap”, tangkap habis tidak tahu untuk siapa. Sebelum mula tangkap, dahulu jawab tiga soalan:

\begin{enumerate}
  \item \textbf{Guna stylus untuk nota tulisan tangan, lukisan, edit foto.}- Tangkap untuk anak tengok, tangkap untuk rakan sebaya tengok, tangkap untuk warga emas tengok, gaya kandungan berbeza.
  \item \textbf{AI (Kecerdasan Buatan) telah menjadi sebahagian besar telefon pintar moden. Bab ini membincangkan ciri AI yang ada pada telefon dan cara menggunakannya.}I. Asas AI telefon
  \item \textbf{1. NPU (Neural Processing Unit)}- Douyin/Kuaishou (hiburan umum), Xiaohongshu (gaya hidup), Bilibili (kandungan mendalam), Video Account (bulatan sosial). Pilih satu pangkalan utama mendalami, jangan mula-mula semua platform
\end{enumerate}

\booksection{Tetapan asas rakaman telefon — 5 perkara wajib sebelum merakam}

\begin{enumerate}
  \item \textbf{2. AI pada peranti vs AI awan}: Ini adalah standard platform video pendek. Fail 4K terlalu besar muat naik perlahan, 60fps di platform sosial tiada kelebihan sebaliknya fail lebih besar.
  \item \textbf{AI awan: Diproses di pelayan, lebih berkuasa, tetapi perlu sambungan internet dan kurang privasi.}3. Model bahasa besar (LLM) pada telefon
  \item \textbf{Telefon flagship 2024-2026 mula menjalankan LM (seperti Gemini Nano, Apple Intelligence) pada peranti untuk tugas seperti ringkasan teks, janaan imej, penterjemahan.}II. Ciri AI telefon
  \item \textbf{1. Fotografi AI}- Hantar Douyin/Kuaishou/Video Account/Xiaohongshu\textbf{HDR AI: AI proses berbilang dedahan untuk gambar HDR yang lebih baik.}(9:16), hantar Bilibili/YouTube\textbf{2. AI dalam video}Penstabilan video: AI stabilkan video walaupun berjalan.
  \item \textbf{Pengasingan subjek: AI pisahkan subjek dari latar untuk efek bokeh atau ganti latar.}: Tetapan > Kamera > Garis grid. Asas komposisi rule of thirds, bantu anda cepat cari pusat graviti gambar.
\end{enumerate}

\booksection{Lima teknik pergerakan kamera — Jadikan imej “hidup”}

Masalah terbesar pemula tangkap video adalah gambar terlalu “senyap” - sepanjang masa berdiri tidak gerak tangkap. Belajar lima pergerakan kamera asas berikut, gambar serta-merta ada rasa profesional:

\begin{itemize}
  \item \textbf{Google Assistant: Pembantu Google, sokong perbualan semula jadi, tugas kompleks.}Xiaomi Xiaoai: Pembantu suara domestik, sokong kawalan peranti pintar.
  \item \textbf{4. AI dalam papan kekunci}Ramalan pintar: AI ramal kata seterusnya berdasarkan konteks.
  \item \textbf{Pembetulan auto: AI betulkan ejaan dan tatabahasa.}Penterjemahan: AI terjemah teks semasa taip.
  \item \textbf{5. AI dalam App}Ringkasan teks: AI ringkaskan artikel, e-mel, dokumen panjang.
  \item \textbf{Pengecaman imej: AI kenalpasti objek, tempat, orang dalam foto.}Janaan imej: AI jana imej berdasarkan teks (seperti DALL-E, Midjourney).
\end{itemize}

\textbf{6. AI kesihatan}Analisis tidur: AI analisis kualiti tidur dan beri cadangan.

\booksection{Asas penyuntingan — CapCut sudah cukup}

Degupan jantung tidak normal: AI kesan degupan jantung tidak normal dan maklum.

\begin{enumerate}
  \item \textbf{7. AI keselamatan}: Buka JianYing > Mula cipta > Pilih klip video yang telah dipotong import.
  \item \textbf{Pengesanan phishing: AI kesan pautan phishing dalam mesej dan e-mel.}Pengesanan malware: AI kesan tingkah laku malware dan maklum.
  \item \textbf{8. Apple Intelligence}: Antara dua segmen video ketik butang "transisi", pilih satu kesan transisi mudah. Cadangan pemula guna "dissolve" atau "fade to black", jangan guna transisi 3D berbelit-belit.
  \item \textbf{Writing Tools: Alat penulisan AI untuk ringkasan, pembetulan, perubahan nada.}Image Playground: Jana imej berdasarkan deskripsi teks.
  \item \textbf{Genmoji: Jana emoji tersuai berdasarkan deskripsi teks.}: Pustaka muzik JianYing ada banyak muzik bebas hak cipta. Pilih satu BGM rentak dan kandungan sepadan, volum laras ke kira-kira 1/3 audio asal video.
  \item \textbf{ChatGPT integrasi: Sokong ChatGPT untuk soalan kompleks.}9. Galaxy AI (Samsung)
\end{enumerate}

\booksection{Kepentingan muka depan — 70\% orang putuskan tonton atau tidak hanya dari muka depan}

Cover video pendek adalah “muka”, tentukan orang mahu klik masuk atau tidak. Beberapa perkara buat cover:

\begin{itemize}
  \item \textbf{Live Translate: Terjemah perbualan telefon real-time.}: Tulis kandungan teras video dengan fon besar, contoh "3 minit belajar pembaikan foto telefon". Fon mesti besar, warna terang, ada kontras dengan latar belakang.
  \item \textbf{Photo Assist: Edit foto dengan AI, termasuk generative fill.}10. Google Gemini
  \item \textbf{Gemini adalah model AI Google, disepadukan dalam Android:}Gemini Nano: Model pada peranti untuk tugas asas, privasi.
  \item \textbf{Gemini Pro/Ultra: Model awan untuk tugas kompleks.}Magic Compose: AI jana draf mesej dengan pelbagai nada.
\end{itemize}

\booksection{Tiga prinsip penciptaan kandungan}

\begin{itemize}
  \item \textbf{11. AI domestik (China)}: Platform video pendek, penonton 3 saat tentukan langkau atau terus tengok. 3 saat pertama mesti beri "pengait" - satu soalan, satu misteri, satu gambar mengejutkan.
  \item \textbf{Xiaomi: Super Xiaoai, sokong perbualan semula jadi, kawalan peranti pintar.}: Kesilapan paling mudah pemula adalah satu video terlalu banyak kandungan. Satu video hanya selesaikan satu masalah, terangkan mendalam lebih penting daripada terangkan semua.
  \item \textbf{III. Cara guna AI telefon}: Algoritma cadangan perlu "pelihara akaun". 20 video pertama hampir tiada jumlah tontonan, ini adalah normal. Berterusan kemas kini 2-3 video seminggu, biar algoritma kenal arah kandungan anda, trafik akan perlahan-lahan naik.
\end{itemize}

\begin{tcbpractice}{Praktikal: Tangkap satu video pendek 30 saat}
\begin{enumerate}
  \item 2. Guna AI untuk produktiviti
  \item Ringkasan: Guna AI untuk ringkaskan artikel panjang, e-mel, dokumen.
  \item Penterjemahan: Guna AI untuk terjemah teks, suara, imej.
  \item Penulisan: Guna AI untuk bantu tulis, betulkan, ubah nada teks.
  \item 3. Guna AI untuk fotografi
\end{enumerate}

\end{tcbpractice}

\begin{tcbsummary}{Windows 11: Tetapan Sistem Papar Pelbagai paparan Sambung ke paparan wayarles.}
\begin{itemize}
  \item Biarkan AI pilih scene: Buka kamera, biarkan AI kenalpasti scene dan optimumkan.
  \item Guna mod malam AI: Untuk foto malam yang lebih jelas.
  \item Edit foto dengan AI: Guna alat edit AI untuk buang objek, ganti latar, tambah efek.
  \item 4. Guna AI untuk kesihatan
  \item 3 saat pertama tentukan hidup mati, satu video hanya cakap satu perkara, konsisten lebih penting daripada bakat
\end{itemize}

\end{tcbsummary}
